\chapter{RSA Cryptosystem}

\section{What is RSA?}
RSA is one of the most famous public-key cryptosystems. It lets people
communicate securely without first sharing a secret key.

The core idea is simple:
\begin{itemize}
  \item A \textbf{public key} is shared with everyone and is used to encrypt.
  \item A \textbf{private key} is kept secret and is used to decrypt.
\end{itemize}

RSA is also used for digital signatures, but in this chapter we focus on
encryption/decryption and the mathematical principle behind why it works.

\subsection*{Alice, Bob, and Eve}
Think of three people:
\begin{itemize}
  \item \textbf{Alice}: wants to receive a secret message.
  \item \textbf{Bob}: wants to send the message.
  \item \textbf{Eve}: can eavesdrop on every message sent over the network.
\end{itemize}

Alice publishes her \textbf{public key}. Bob uses that public key to encrypt.
Eve can see the public key and the ciphertext, but cannot recover the plaintext
without Alice's private key. Alice then uses her private key to decrypt.

So the story is:
\begin{enumerate}
  \item Alice generates keys and publishes \((n,e)\).
  \item Bob computes \(c\equiv m^e\pmod n\) and sends \(c\).
  \item Eve sees \(c\), but should not be able to get \(m\).
  \item Alice computes \(m\equiv c^d\pmod n\) and recovers the message.
\end{enumerate}

\section{Key Generation}
Choose two large prime numbers \(p\) and \(q\), then define
\[
  n = pq.
\]
Compute
\[
  \varphi(n) = (p-1)(q-1).
\]
Choose an integer \(e\) such that
\[
  1 < e < \varphi(n), \qquad \gcd(e,\varphi(n))=1.
\]
Then choose \(d\) such that
\[
  ed \equiv 1 \pmod{\varphi(n)}.
\]
The public key is \((n,e)\), and the private key is \((n,d)\).

\section{Encryption and Decryption}
Represent a message as an integer \(m\) with \(0 \le m < n\).

\textbf{Encryption:}
\[
  c \equiv m^e \pmod n.
\]

\textbf{Decryption:}
\[
  m \equiv c^d \pmod n.
\]

In words: raise to \(e\) to encrypt, raise to \(d\) to decrypt.

In the Alice-Bob-Eve picture:
\begin{itemize}
  \item Bob uses Alice's public key \((n,e)\) to compute \(c\).
  \item Eve can read \(c\), but she does not know Alice's \(d\).
  \item Alice uses \(d\) to recover \(m\).
\end{itemize}

\section{Small Numerical Example}
For teaching, we use small numbers (in practice, primes are much larger):
\[
  p=61,\quad q=53.
\]
Then
\[
  n = pq = 3233,\qquad \varphi(n)=(61-1)(53-1)=3120.
\]
Pick
\[
  e=17,\qquad \gcd(17,3120)=1.
\]
Find \(d\) from \(ed\equiv 1 \pmod{3120}\), which gives
\[
  d=2753
\]
because \(17\cdot 2753 = 46801 = 3120\cdot 15 + 1\).

So:
\begin{itemize}
  \item Public key: \((n,e)=(3233,17)\)
  \item Private key: \((n,d)=(3233,2753)\)
\end{itemize}

Encrypt message \(m=65\):
\[
  c \equiv 65^{17} \pmod{3233} = 2790.
\]

Decrypt:
\[
  m \equiv 2790^{2753} \pmod{3233} = 65.
\]
So the original message is recovered.

\section{Why RSA Works}
The proof of correctness is based on Fermat's little theorem:
if \(p\) is prime and \(p\nmid a\), then
\[
  a^{p-1}\equiv 1 \pmod p.
\]

We want to show
\[
  (m^e)^d \equiv m \pmod{pq}
\]
for every integer \(m\), where \(p\) and \(q\) are distinct primes, and \(e,d\)
satisfy
\[
  ed \equiv 1 \pmod{\varphi(pq)}.
\]


\begin{enumerate}
  \item gcd(m, pq) \neq 1, 

  Both p and q are primes. m must be multiple of p or q. let's assume m is multiple of p. m = pk. m is multiple
  \[
    m^{ed}
    =(pk)^{ed}
    \equiv 0 
    \equiv m \pmod p
  \]

  \[
    m^{ed}
    =m^{ed-1}m
    =m^{k(q-1)}m
    \equiv m \pmod q
  \]

  We have therefore
  \[
    m^{ed}=qx+m
  \]
  Since $m^{ed}$ is multiple of p, m is multiple of p, we have x is multiple of p.
  \[
    m^{ed}=pq(x/p)+m
    \equiv m \pmod {pq}
  \]

  \item gcd(m, pq) = 1, 
  We can use Euler's theorem to prove directly.
  \[
    m^{ed}
    =m^{ed-1}m
    =m^{k\varphi(pq)}m
    \equiv 1^k m
    \equiv m
    \pmod {pq}
  \]
\end{enumerate}

Therefore \(m^{ed}\equiv m\pmod p\) and \(m^{ed}\equiv m\pmod q\), so
\[
  (m^e)^d \equiv m \pmod{pq}.
\]
This proves correctness for every integer \(m\).

\section{Modular Exponentiation}
At this point, a practical question appears:
\[
  m \equiv 2790^{2753} \pmod{3233}.
\]
How can we compute this without multiplying \(2790\) by itself 2753 times?

The naive method is too slow for real RSA. In practice, exponents are huge
(hundreds or thousands of bits), so we must use a smarter method.

The standard trick is \textbf{binary exponentiation} (also called
square-and-multiply). It uses:
\[
  a^b=
  \begin{cases}
    (a^{b/2})^2, & b \text{ even},\\[2mm]
    a\,(a^{(b-1)/2})^2, & b \text{ odd}.
  \end{cases}
\]
So instead of reducing the exponent by \(1\) each step, we roughly halve it
each step. That is why the complexity is \(O(\log b)\), not \(O(b)\).

We also reduce modulo \(n\) at every multiplication:
\[
  (xy)\bmod n = \bigl((x\bmod n)(y\bmod n)\bigr)\bmod n.
\]
This keeps numbers small during the whole computation.

Here is a clean recursive implementation:
\begin{lstlisting}[language=Python]
def fast_expo_modulo(base: int, exp: int, modulo: int):
    if exp==1:
        return base%modulo
    if exp ==0:
        return 1
    if exp % 2 == 1:
        return (fast_expo_modulo(base, (exp-1)/2, modulo) %modulo* 
            fast_expo_modulo(base, (exp-1)/2, modulo) %modulo 
            * (base%modulo))%modulo
    else:
        return (fast_expo_modulo(base, exp/2, modulo) %modulo* 
            fast_expo_modulo(base, exp/2, modulo) %modulo) %modulo
\end{lstlisting}

\section{Computing the Private Key with Extended Euclid}
In RSA key generation, after choosing \(e\), we must compute the private key
exponent \(d\) from
\[
  ed \equiv 1 \pmod{\varphi(n)}
\]
(or modulo \(\lambda(n)\) in many implementations). So \(d\) is the
\textbf{modular multiplicative inverse} of \(e\).

The Extended Euclidean Algorithm gives exactly this inverse. Since
\(\gcd(e,\varphi(n))=1\), it finds integers \(x,y\) such that
\[
  ex+\varphi(n)y=1.
\]
Taking both sides modulo \(\varphi(n)\):
\[
  ex \equiv 1 \pmod{\varphi(n)}.
\]
Therefore \(x\) is an inverse of \(e\), so we can take
\[
  d \equiv x \pmod{\varphi(n)}.
\]
If \(x<0\), convert it to a positive representative by adding \(\varphi(n)\).
This is the standard method used in RSA implementations \cite{wiki:rsa-cryptosystem,wiki:modular-multiplicative-inverse}.

\subsection*{How the algorithm works (and why)}
The ordinary Euclidean algorithm repeatedly applies division with remainder:
\[
a=bq+r,\qquad 0\le r<b.
\]
Then \(\gcd(a,b)=\gcd(b,r)\). Repeating this step makes numbers smaller until
the remainder is \(0\), and the last nonzero remainder is the gcd.

The \emph{extended} Euclidean algorithm keeps extra bookkeeping: at each step,
it remembers how the current remainder is written as a linear combination of
the original two numbers. So every remainder has the form
\[
r_i = s_i a + t_i b.
\]
When the gcd is \(1\), this finally gives
\[
1 = sa+tb.
\]
This identity is called a B\'{e}zout identity.

Now apply it to RSA with \(a=e\) and \(b=\varphi(n)\):
\[
1 = se + t\varphi(n).
\]
Reduce modulo \(\varphi(n)\):
\[
se \equiv 1 \pmod{\varphi(n)}.
\]
So \(s\) is exactly \(e^{-1}\bmod \varphi(n)\), i.e. the private exponent
\(d\) (up to adding multiples of \(\varphi(n)\)).

\subsection*{Worked Example}
Use the chapter's RSA toy parameters:
\[
  e=17,\qquad \varphi(n)=3120.
\]
Euclidean steps:
\[
3120=183\cdot17+9,\quad
17=1\cdot9+8,\quad
9=1\cdot8+1.
\]
Back-substitute:
\[
\begin{aligned}
1&=9-1\cdot8\\
 &=9-(17-1\cdot9)=2\cdot9-17\\
 &=2(3120-183\cdot17)-17\\
 &=2\cdot3120-367\cdot17.
\end{aligned}
\]
So
\[
17(-367)+3120(2)=1.
\]
Hence
\[
17(-367)\equiv1\pmod{3120},
\]
and the private exponent is
\[
d\equiv -367\equiv 2753\pmod{3120}.
\]
This is exactly the \(d\) used earlier in the chapter.

\section{Practical Notes}
\begin{itemize}
  \item Real RSA uses very large primes (typically 2048-bit keys or larger).
  \item RSA security relies on the difficulty of factoring large integers.
\end{itemize}

\section*{Reference}
This chapter follows the standard RSA description in \cite{wiki:rsa-cryptosystem}.
